% SYSTEM IMPLEMENTATION AND VALIDATION
% Aligned with the Tawakad Flutter codebase (lib/, android/, ios/).
% Compile as part of your main thesis; adjust packages (babel/polyglossia) if needed.

\chapter{SYSTEM IMPLEMENTATION AND VALIDATION}

\section{Introduction}

This chapter presents the implementation and validation processes of the Tawakad
application. It describes how the system requirements introduced in previous chapters
were transformed into functional system components through an iterative Agile-Scrum
development process.

\vspace{0.3cm}

The implementation follows a \textbf{cross-platform Flutter architecture}: a shared
Dart implementation under \texttt{lib/} targets Android and iOS, with small amounts
of native code where platform capabilities require it (notably the iOS calendar
integration). Cloud services are provided through \textbf{Firebase Authentication} and
\textbf{Cloud Firestore} for user-centric persistence. Intelligent recommendations are
\textbf{specified in the product experience} (including onboarding messaging and
stored questionnaire answers); the current repository focuses on \textbf{data
foundations and client UX}, while full machine-learning deployment is treated as a
\textbf{planned extension} rather than a shipped training and inference pipeline in
this codebase.

\vspace{0.3cm}

The chapter also explains the implementation details of the user interface, client-side
state management, Firebase integration, and platform-specific modules. Validation and
testing procedures are discussed to evaluate functionality, usability, and reliability
within the scope of the implemented prototype.

\vspace{0.3cm}

The implementation activities were organized into multiple development sprints aligned
with the sprint planning methodology and system requirements defined in the Software
Requirements Specification (SRS).

% ============================================================
\section{Implementation by Sprints}

% ------------------------------------------------------------
\subsection{Sprint 1: Authentication Module}

This sprint focused on implementing the user authentication system, including login,
registration, email verification, password recovery, and secure linkage to cloud user
records. In addition to email and password authentication, \textbf{Google Sign-In} was
integrated as an alternative identity provider.

% ------------------------------------------------------------
\subsubsection{Frontend Implementation}

The frontend implementation included:

\begin{itemize}
    \item Designing login and registration user interfaces (including Arabic UX copy),
    \item Responsive authentication and onboarding-related screens,
    \item Client-side form validation and input checking (including password-strength rules on sign-up),
    \item Error message handling for authentication failures,
    \item Password visibility controls where applicable,
    \item Navigation between authentication pages (sign-in, sign-up, verification, forgot password),
    \item Integration with Firebase Authentication APIs (email/password and Google credential flows).
\end{itemize}

% ------------------------------------------------------------
\subsubsection{Backend Implementation}

The backend implementation included:

\begin{itemize}
    \item Firebase Authentication integration for account lifecycle operations,
    \item User account creation and management (including Google-based accounts),
    \item Email verification workflow (verification email dispatch and enforcement before sign-in for email/password users),
    \item Password reset services via Firebase,
    \item Authentication state as managed by the Firebase client SDK (session continuity on the device),
    \item Associated Firestore user document creation to store onboarding answers and password-history metadata (hashed) for policy support.
\end{itemize}

% ============================================================
\subsection{Sprint 2: Home Dashboard and List Management Module}

This sprint focused on implementing the main dashboard experience for \textbf{pack lists}
(items the user intends to bring or check), including creation and editing flows,
filtering, favorites, and in-app search. List state is managed on the \textbf{client}
during this implementation phase.

% ------------------------------------------------------------
\subsubsection{Frontend Implementation}

The frontend implementation included:

\begin{itemize}
    \item Home dashboard UI design and layout,
    \item Pack list card rendering and list-centric interactions,
    \item Dynamic filtering controls (for example, today, favorites, and all lists),
    \item Favorite indicators and toggling,
    \item Multi-step creation and editing flows (metadata such as date, time, repeat options, and sharing flags),
    \item State management for list updates using a provider pattern,
    \item Search-driven list filtering from the application shell,
    \item Bottom navigation shell with Home, Calendar, and a placeholder tab reserved for future scanning features.
\end{itemize}

% ------------------------------------------------------------
\subsubsection{Backend Implementation}

The backend implementation included:

\begin{itemize}
    \item Firebase project configuration for the mobile clients (core initialization),
    \item \textbf{Firestore integration used for user profile documents} (not for per-list document CRUD in the current codebase),
    \item In-memory list persistence within the running app session (prototype scope),
    \item Client-side scheduling-related fields captured on list entities (date/time/repeat), without a separate push-notification scheduling service in the current dependency set,
    \item Data modeling for lists and items suitable for future cloud persistence if extended.
\end{itemize}

% ============================================================
\subsection{Sprint 3: Application Shell, Global Presentation, and Navigation}

This sprint focused on implementing the \textbf{application shell} and global presentation
defaults that affect the whole product, rather than a full standalone settings module.
The codebase configures theming and localization at the application root, and provides
primary navigation between major areas.

% ------------------------------------------------------------
\subsubsection{Frontend Implementation}

The frontend implementation included:

\begin{itemize}
    \item Application shell with animated bottom navigation and integrated search affordances,
    \item Global light and dark themes with system-driven theme mode selection where configured,
    \item Localization scaffolding (for example, Arabic as the active locale with additional locales registered for future expansion),
    \item Consistent typography and shared UI components across modules,
    \item Navigation between Home, Calendar, and reserved tabs,
    \item Placeholder screens where platform-specific or future functionality is not yet implemented (for example, non-iOS calendar presentation and scanning tab).
\end{itemize}

% ------------------------------------------------------------
\subsubsection{Backend Implementation}

The backend implementation included:

\begin{itemize}
    \item Continued reliance on Firebase-backed identity for associating future per-user cloud data with authenticated users,
    \item Firestore-backed user document fields already created at registration time (onboarding answers) as a foundation for later personalization features,
    \item No dedicated Firestore settings document model in the current prototype scope (settings persistence can be added as a future sprint).
\end{itemize}

% ============================================================
\subsection{Sprint 4: Intelligent Recommendations: Data Foundations and Roadmap}

This sprint focused on \textbf{establishing data foundations and product alignment} for
an intelligent recommendation experience, while clearly separating \textbf{what is
implemented now} from \textbf{planned machine-learning components}.

% ------------------------------------------------------------
\subsubsection{Data Collection}

\textbf{Implemented in the current prototype:} onboarding questionnaire answers are
captured in the client flow and persisted in the user's Firestore document at registration,
alongside basic account metadata.

\vspace{0.2cm}

\textbf{Planned or partially available inputs for future recommendation logic:} reminder
and list usage history once cloud persistence is extended, richer preference profiles,
calendar context (the iOS implementation embeds a native calendar view for interaction
with device calendars), environmental context such as weather, and longitudinal usage
patterns. These sources are part of the system vision described in earlier requirements,
but not all are fully wired into an automated training dataset in the current repository.

% ------------------------------------------------------------
\subsubsection{Data Preprocessing and Cleaning}

For the data that is currently stored (notably onboarding answers), lightweight validation
is enforced at input time in the client application. A full preprocessing pipeline for
machine learning (normalization, deduplication, labeling, and feature engineering across
heterogeneous signals) is \textbf{planned} as part of a future analytics and model-training
workflow, rather than implemented as a batch processing service in this codebase snapshot.

% ------------------------------------------------------------
\subsubsection{Model Comparison}

Model comparison (for example, classical baselines versus contextual models, on-device
constraints versus cloud inference) is part of the \textbf{planned evaluation strategy}
for the recommendation subsystem. At the prototype stage documented here, the mobile
client does not include an integrated training or batch evaluation harness; comparative
experiments would be conducted in an offline research environment before model selection
for deployment.

% ------------------------------------------------------------
\subsubsection{Selected Model}

No single deployed ML model is claimed as the final production recommender within this
implementation chapter. Instead, the selection criteria are recorded as requirements:
\textbf{contextual adaptability}, \textbf{latency acceptable for mobile UX}, and
\textbf{maintainability} (including privacy and updateability). The eventual model choice
should be justified against those criteria once training data and serving infrastructure
are in place.

% ------------------------------------------------------------
\subsubsection{Evaluation Metrics}

When a recommendation model is integrated, evaluation is expected to include standard
classification-style metrics where applicable (accuracy, precision, recall, F1-score)
plus product-facing measures such as \textbf{recommendation relevance} and
\textbf{user satisfaction} in usability studies. For the prototype without an integrated
model server, AI-specific quantitative validation is deferred; the validation section below
focuses on implemented flows.

% ============================================================
\section{Validation and Testing}

This section describes the testing and validation procedures used to evaluate the
correctness, reliability, usability, and performance of the Tawakad application within
the scope of the implemented features.

% ------------------------------------------------------------
\subsection{UI Validation}

UI validation focused on evaluating:

\begin{itemize}
    \item Visual consistency across authentication, onboarding, and home flows,
    \item Layout behavior across common phone sizes (Android and iOS),
    \item Basic accessibility considerations (contrast, tap targets) as applicable,
    \item Navigation flow between shell tabs and authentication routes,
    \item Readability of Arabic UI strings and hierarchy of primary actions.
\end{itemize}

% ------------------------------------------------------------
\subsection{Backend and Firebase Validation}

Backend validation focused on:

\begin{itemize}
    \item Email/password authentication workflows, including verification and password reset,
    \item Google Sign-In credential exchange and Firebase user creation or linkage,
    \item Firestore write and read behavior for the user document created at registration,
    \item Client-side error handling for common failure modes (network, invalid credentials, cancelled Google sign-in),
    \item Email-based account recovery paths provided by Firebase.
\end{itemize}

% ------------------------------------------------------------
\subsection{AI and Recommendation Readiness}

Validation for the intelligent recommendation subsystem is currently \textbf{requirements-level
and architectural}: onboarding data is persisted to support future personalization, and
the UX communicates the intended assistant behavior. \textbf{Quantitative model validation}
(offline metrics, A/B tests, latency under load) is planned for a later milestone when a
model and serving approach are integrated.

% ------------------------------------------------------------
\subsection{System Testing Procedures}

The following testing methods were applied or are applicable to the prototype scope:

\begin{itemize}
    \item Functional testing of authentication, list creation/editing, filtering, favorites, and search,
    \item Integration testing across Flutter UI, Firebase SDKs, and (on iOS) native calendar platform views,
    \item Usability testing for primary task completion (sign up/in, create a list, filter, favorite),
    \item Lightweight performance checks for typical navigation and list sizes on mobile devices,
    \item Compatibility testing across Android and iOS builds from the same \texttt{lib/} codebase, noting platform differences (for example, calendar tab behavior),
    \item User acceptance testing (UAT) with representative users once stable builds are distributed.
\end{itemize}

% ============================================================
\section{Summary}

This chapter presented the implementation and validation phases of the Tawakad
application. The system was developed incrementally through Agile-based development sprints
covering authentication (including Google Sign-In), pack list management with client-side
state in the prototype scope, an application shell with navigation and global presentation
defaults, and the \textbf{foundations and roadmap} for contextual recommendations backed
by onboarding data stored in Firestore. Validation emphasized implemented Firebase workflows,
UI consistency, and cross-platform behavior, while deferring full AI model evaluation until
a recommender is integrated into the product.
